\chapter{Introduction}\label{ch:intro}

\section{Topic and motivation}

Generative diffusion models have become the dominant paradigm for synthesising
high-dimensional data: they power state-of-the-art systems for images, audio,
video, and molecular structures \citep{sohl2015deep, ho2020denoising,
song2021scorebased, karras2022elucidating}. Their central object is the
\emph{score function} --- the gradient of the log-density of progressively
noised data,
\begin{equation}
  \score(x, t) \;=\; \nabla_x \log p_t(x),
  \label{eq:intro-score}
\end{equation}
which is all a model needs to know in order to reverse a noising process and
turn pure noise back into data \citep{anderson1982reverse}. In practice the
score is approximated by a neural network trained by denoising score matching
\citep{hyvarinen2005estimation, vincent2011connection}; the network is a black
box, and the structure of what it has to learn remains poorly understood.

A complementary research programme, initiated in the statistical-physics
community, replaces the black box with \emph{exactly solvable models}: data
distributions simple enough that the score can be computed in closed form, yet
rich enough to exhibit the phenomena observed in real systems --- speciation,
memorisation, and dynamical phase transitions along the generation trajectory
\citep{biroli2023generative, biroli2024dynamical, raya2023spontaneous,
ghio2024sampling}. Exact scores act as a microscope: every regime, transition
time, and failure mode can be traced back to explicit formulas.

This thesis extends the exactly solvable programme in a direction that is
natural for data with \emph{temporal or sequential structure}: the data are not
a cloud of independent points but a \emph{trajectory} $a_0, a_1, \dots,
a_{K-1}$ of a Markov process --- the prototype of a video, a time series, or
any dynamic object. Each frame of the trajectory is corrupted independently by
an Ornstein--Uhlenbeck (OU) channel, and the object of study is the
\emph{joint} score of the whole noisy sequence. Two classical bodies of theory
then meet the modern one. First, the joint score of a Markov chain turns out
to be governed by the algebra of structured matrices --- Toeplitz covariances,
tridiagonal precisions, and their spectral properties. Second, computing the
score is an exercise in Bayesian inference on a tree-shaped factor graph, the
home turf of \emph{belief propagation} (BP) and its Gaussian projection,
\emph{approximate message passing} (AMP) \citep{pearl1988probabilistic,
mezard2009information, donoho2009message}. The thesis develops both routes in
full detail, proves where they coincide, and measures precisely what is lost
when they are forced apart.

\section{Positioning and gap}

The literature this work builds on can be organised in three strands.

\paragraph{Statistical physics of diffusion models.} Closed-form analyses of
diffusion generation exist for Gaussian mixtures, random energy models, and
high-dimensional asymptotics \citep{biroli2023generative, biroli2024dynamical,
ghio2024sampling}. These works treat the data as unstructured collections of
samples; the time axis they study is the \emph{diffusion} time. The internal
time of a \emph{sequence} --- the fact that a video frame is correlated with
its neighbours through a dynamical rule --- is absent from these analyses.

\paragraph{Inference on chains.} Exact posterior inference on Gauss--Markov
chains is classical: the Kalman filter and smoother
\citep{kalman1960new, rauch1965maximum}, hidden Markov models
\citep{rabiner1989tutorial}, and their unifying formulation as sum--product
message passing on factor graphs \citep{kschischang2001factor,
mezard2009information}. What this literature does not ask is the question that
matters for generative modelling: \emph{how does the exact posterior --- and
therefore the exact score --- deform as a function of the diffusion time},
and what does that deformation imply for the architecture of a network that
must learn it?

\paragraph{Message passing with continuous variables.} BP is exact on trees
for discrete variables, where messages are finite vectors of probabilities.
For continuous variables the messages become probability \emph{densities} ---
infinite-dimensional objects --- and practical algorithms must close them
within a parametric family, almost always the Gaussian one; this is the step
that turns BP into AMP, known in physics as the
Thouless--Anderson--Palmer (TAP) approach \citep{thouless1977solution,
donoho2009message, zdeborova2016statistical}. The classical justification of
the Gaussian closure is a central-limit argument valid when every node
receives many messages. On a chain each node receives exactly two: the
central-limit argument fails by construction, and the quality of the Gaussian
closure on chain-structured diffusion posteriors was, to the best of our
knowledge, unquantified.

The gap, in one sentence: \emph{there is no end-to-end, exactly solvable
account of the joint diffusion score of a structured sequence --- its matrix
structure, its message-passing computation, and the price of the Gaussian
closure when the sequence is not Gaussian.}

\section{Research questions}

The thesis addresses four questions.

\begin{itemize}
  \item[\textbf{RQ1.}] What is the exact joint score of a Markov chain
  corrupted by an OU channel, and how does its structure --- concretely, the
  posterior precision matrix --- evolve along the diffusion time, from the
  sparse Markov regime at $t \to 0$ to the isotropic regime at $t \to \infty$?

  \item[\textbf{RQ2.}] Can the joint score be computed \emph{locally}, by
  message passing on the chain's factor graph, rather than globally by matrix
  inversion? In the Gaussian case, how does belief propagation relate to the
  classical Kalman smoother, and does it reproduce the closed-form score
  exactly?

  \item[\textbf{RQ3.}] When the messages are forced into a Gaussian family ---
  the AMP/TAP closure --- what survives and what breaks? On the Gaussian
  chain, is the closure exact; and on a non-Gaussian (Laplace-innovation)
  chain, how large is the closure error as a function of the diffusion time?

  \item[\textbf{RQ4.}] What do the exact results imply for \emph{learned}
  scores: when is a local (banded, convolution-like) score architecture
  sufficient, and what is the exact law governing the error of local
  truncations?
\end{itemize}

\section{Contributions}

The main contributions, each developed in a dedicated chapter, are the
following.

\begin{enumerate}
  \item \textbf{A complete closed-form solution of the Gaussian benchmark}
  (Chapter~\ref{ch:gaussian}). For the stationary AR(1) chain under OU
  corruption we derive the joint score, diagonalise it in closed form, and
  characterise the full lifecycle of the posterior precision matrix: a
  corrected \emph{band-fill law} for the small-$t$ growth of its
  off-diagonals, and exact bulk formulas for its inverse --- variance,
  geometric decay rate $q(\alpha, t)$ of posterior correlations, and their
  limits. The band-fill law corrects a transcription error in an earlier
  working note; the correction is validated numerically.

  \item \textbf{An equation-by-equation message-passing construction of the
  score} (Chapter~\ref{ch:bp}). We build the factor graph of the noisy chain,
  specialise the sum--product equations to it, prove that Gaussian messages
  are closed under the updates (two elementary facts suffice), and show that
  the resulting algorithm --- Gaussian BP, algebraically identical to a
  Kalman smoother --- reproduces the matrix score to machine precision at
  cost $O(K)$ instead of $O(K^3)$.

  \item \textbf{An exact phase boundary for the AMP closure}
  (Chapter~\ref{ch:bp}). On the chain, BP and AMP produce the \emph{same
  score} but \emph{different posterior variances}. We show the difference is
  a single factor of two in a closure discriminant ($J_d^2 - 4\beta^2$ for BP
  against $J_d^2 - 8\beta^2$ for AMP), prove that the BP fixed point exists
  for every coupling and diffusion time, and derive the exact AMP breakdown
  time $t_c(\alpha)$ together with the critical coupling
  $\alpha_c = \sqrt{2} - 1$ below which AMP never breaks down.

  \item \textbf{A quantitative audit of the Gaussian closure beyond
  Gaussianity} (Chapters~\ref{ch:laplace} and~\ref{ch:bp}). For
  Laplace-innovation chains --- where exact inference is still available
  through deterministic grid quadrature on the same factor graph --- we
  measure the score error of Gaussian-projected BP: it peaks around
  $20\%$ (relative $L^2$) at small diffusion time and decays below $1\%$ as
  the noisy marginals Gaussianise, while the score \emph{direction} remains
  accurate throughout. We also show that at small $t$ \emph{truncating the
  inference} (local BP sweeps) is markedly more accurate than \emph{truncating
  the score matrix} (banding), quantifying the advantage of local message
  passing over local linear algebra.

  \item \textbf{A reproducible computational companion.} All results are
  backed by open, dependency-light Python code with deterministic numerical
  audits ($72/72$ and $55/55$ independent checks, re-executed at build
  time); the code and the
  figures it generates are part of the thesis repository and are listed in
  Appendix~\ref{app:code}.
\end{enumerate}

\section{Methodology}

The methodology is deliberately old-fashioned: \emph{no neural networks, no
black boxes --- closed forms first, numerics as audit}. Every claim is either
proved by explicit algebra (the derivations are spelled out step by step, with
the longer algebraic passages isolated in numbered \emph{toolboxes}) or
verified by deterministic numerical checks against brute-force ground truth,
in the spirit of the exactly solvable programme of
\citet{biroli2023generative}. Where an approximation is introduced (the
Gaussian closure of BP messages, local truncations), its error is measured
against an exact reference computed by an independent method (matrix algebra
or dense grid quadrature). The empirical setting, the audit protocol, and the
software packages are described alongside the results and in
Appendix~\ref{app:repro}.

\section{Structure of the thesis}

The thesis is organised in two movements: a foundations movement
(Chapters~\ref{ch:statmech}--\ref{ch:diffusion}) that builds, from first
principles, every tool the research movement needs, and a research movement
(Chapters~\ref{ch:gaussian}--\ref{ch:bp}) that contains the original
contributions.

Chapter~\ref{ch:statmech} develops statistical mechanics from Hamiltonian
mechanics: phase space, Liouville's theorem, the Boltzmann--Gibbs
distribution, free energy, and the Ising paradigm of complex systems.
Chapter~\ref{ch:stochproc} covers stochastic processes and Markov-chain Monte
Carlo: stationarity, detailed balance, Metropolis--Hastings and Gibbs
sampling, and introduces the AR(1) chain that serves as the running example
of the whole thesis. Chapter~\ref{ch:sde} builds the continuous-time toolkit:
Brownian motion, the It\^o integral and It\^o's lemma, the
Ornstein--Uhlenbeck process --- the corruption channel of diffusion models
--- and the Fokker--Planck equation, closing the loop with the Boltzmann
distribution of Chapter~\ref{ch:statmech}. Chapter~\ref{ch:ebm} treats
energy-based models, Markov random fields, hidden Markov models, and factor
graphs --- the representational language of the research chapters.
Chapter~\ref{ch:diffusion} assembles these pieces into modern diffusion
models: the DDPM construction, the score-SDE formulation, score matching,
Tweedie's identity, and the statistical-physics analyses of generative
diffusion that motivate this work.

Chapter~\ref{ch:gaussian} solves the Gaussian chain exactly (RQ1).
Chapter~\ref{ch:laplace} leaves Gaussianity in the most controlled way
possible --- Laplace innovations --- and shows what the joint score looks like
when it is no longer linear (first half of RQ3). Chapter~\ref{ch:bp}
recomputes everything by message passing: BP, its Gaussian specialisation and
the Kalman connection (RQ2), the AMP closure and its exact phase boundary,
the closure error on non-Gaussian chains (RQ3), and the locality laws that
answer RQ4. Chapter~\ref{ch:conclusions} summarises the findings per research
question, discusses limitations, and lays out the research programme that
follows from them. Three appendices collect Gaussian identities, the code
listings, and the reproducibility protocol.
